\documentclass[10pt,a4paper]{amsart}
\usepackage{amsmath,amssymb,amsthm,graphicx}
% Retained preamble for the Stacks Project chapter snapshot in this directory.
% The chapter's own preamble.tex is a copy of the shared Stacks Project preamble; the block
% below is sliced MECHANICALLY, line for line, from it. Nothing was added, removed or
% redefined; the wrapper class, the xr-hyper/hyperref wiring and the externaldocument list
% are omitted because the assembler supplies the document class.
%
% lines taken from preamble.tex: 3, 5, 9, 11, 12, 156, 157, 158, ... (51 lines total)
%
% package coverage restored from the chapter's own preamble (2026-10-13): xy, xy-2cell

\usepackage[all]{xy}
\xyoption{2cell}
\UseAllTwocells

\newenvironment{reference}{\comment}{\endcomment}
\newenvironment{slogan}{\comment}{\endcomment}
\newenvironment{history}{\comment}{\endcomment}
\theoremstyle{plain}
\newtheorem{theorem}[subsection]{Theorem}
\newtheorem{proposition}[subsection]{Proposition}
\newtheorem{lemma}[subsection]{Lemma}
\theoremstyle{definition}
\newtheorem{definition}[subsection]{Definition}
\newtheorem{example}[subsection]{Example}
\newtheorem{exercise}[subsection]{Exercise}
\newtheorem{situation}[subsection]{Situation}
\theoremstyle{remark}
\newtheorem{remark}[subsection]{Remark}
\newtheorem{remarks}[subsection]{Remarks}
\numberwithin{equation}{subsection}
\def\lim{\mathop{\mathrm{lim}}\nolimits}
\def\colim{\mathop{\mathrm{colim}}\nolimits}
\def\Spec{\mathop{\mathrm{Spec}}}
\def\Hom{\mathop{\mathrm{Hom}}\nolimits}
\def\Ext{\mathop{\mathrm{Ext}}\nolimits}
\def\SheafHom{\mathop{\mathcal{H}\!\mathit{om}}\nolimits}
\def\SheafExt{\mathop{\mathcal{E}\!\mathit{xt}}\nolimits}
\def\Sch{\mathit{Sch}}
\def\Mor{\mathop{\mathrm{Mor}}\nolimits}
\def\Ob{\mathop{\mathrm{Ob}}\nolimits}
\def\Sh{\mathop{\mathit{Sh}}\nolimits}
\def\NL{\mathop{N\!L}\nolimits}
\def\CH{\mathop{\mathrm{CH}}\nolimits}
\def\proetale{{pro\text{-}\acute{e}tale}}
\def\etale{{\acute{e}tale}}
\def\QCoh{\mathit{QCoh}}
\def\Ker{\mathop{\mathrm{Ker}}}
\def\Im{\mathop{\mathrm{Im}}}
\def\Coker{\mathop{\mathrm{Coker}}}
\def\Coim{\mathop{\mathrm{Coim}}}
\DeclareMathSymbol{\boxtimes}{\mathbin}{AMSa}{"02}
\def\QCohstack{\mathcal{QC}\!\mathit{oh}}
\def\Cohstack{\mathcal{C}\!\mathit{oh}}
\def\Spacesstack{\mathcal{S}\!\mathit{paces}}
\def\Quotfunctor{\mathrm{Quot}}
\def\Hilbfunctor{\mathrm{Hilb}}
\def\Curvesstack{\mathcal{C}\!\mathit{urves}}
\def\Polarizedstack{\mathcal{P}\!\mathit{olarized}}
\def\Complexesstack{\mathcal{C}\!\mathit{omplexes}}
\def\Pic{\mathop{\mathrm{Pic}}\nolimits}
\def\Picardstack{\mathcal{P}\!\mathit{ic}}
\def\Picardfunctor{\mathrm{Pic}}
\def\Deformationcategory{\mathcal{D}\!\mathit{ef}}
\usepackage{cleveref}
\crefname{originallemma}{original lemma}{original lemmas}
\Crefname{originallemma}{Original lemma}{Original lemmas}
\makeatletter
\newlabel{algebra-lemma-characterize-dvr}{{00PD}{source}{}{sourcepointer-1}{}}
\newlabel{algebra-lemma-characterize-dvr@cref}{{[originallemma][0][]00PD}{[1][0][]source}{}{}{}}
\makeatother
\crefname{originallemma}{original lemma}{original lemmas}
\Crefname{originallemma}{Original lemma}{Original lemmas}
\makeatletter
\newlabel{divisors-lemma-effective-Cartier-in-points}{{0AG8}{source}{}{sourcepointer-2}{}}
\newlabel{divisors-lemma-effective-Cartier-in-points@cref}{{[originallemma][0][]0AG8}{[1][0][]source}{}{}{}}
\makeatother
\crefname{originallemma}{original lemma}{original lemmas}
\Crefname{originallemma}{Original lemma}{Original lemmas}
\makeatletter
\newlabel{divisors-lemma-fibre-Cartier}{{062Y}{source}{}{sourcepointer-3}{}}
\newlabel{divisors-lemma-fibre-Cartier@cref}{{[originallemma][0][]062Y}{[1][0][]source}{}{}{}}
\makeatother
\crefname{originallemma}{original lemma}{original lemmas}
\Crefname{originallemma}{Original lemma}{Original lemmas}
\makeatletter
\newlabel{more-morphisms-lemma-characterize-finite}{{02LS}{source}{}{sourcepointer-4}{}}
\newlabel{more-morphisms-lemma-characterize-finite@cref}{{[originallemma][0][]02LS}{[1][0][]source}{}{}{}}
\makeatother
\crefname{originallemma}{original lemma}{original lemmas}
\Crefname{originallemma}{Original lemma}{Original lemmas}
\makeatletter
\newlabel{morphisms-lemma-finite-flat}{{02KB}{source}{}{sourcepointer-5}{}}
\newlabel{morphisms-lemma-finite-flat@cref}{{[originallemma][0][]02KB}{[1][0][]source}{}{}{}}
\makeatother
\crefname{originallemma}{original lemma}{original lemmas}
\Crefname{originallemma}{Original lemma}{Original lemmas}
\makeatletter
\newlabel{morphisms-lemma-noetherian-finite-type-finite-presentation}{{01TX}{source}{}{sourcepointer-6}{}}
\newlabel{morphisms-lemma-noetherian-finite-type-finite-presentation@cref}{{[originallemma][0][]01TX}{[1][0][]source}{}{}{}}
\makeatother
\crefname{originallemma}{original lemma}{original lemmas}
\Crefname{originallemma}{Original lemma}{Original lemmas}
\makeatletter
\newlabel{varieties-lemma-smooth-regular}{{056S}{source}{}{sourcepointer-7}{}}
\newlabel{varieties-lemma-smooth-regular@cref}{{[originallemma][0][]056S}{[1][0][]source}{}{}{}}
\makeatother
\title{Divisor conditions on a smooth relative curve}
\author{The Stacks Project Authors}
\date{stacks-project:pic.tex:local-snapshot-20261012; GNU FDL 1.2 or later}
\begin{document}
\maketitle
\noindent\textit{Editorial note.} Statement span is lines 624-642 of sources/stacks-pic/pinned-original-source.tex; proof span is lines 644-675, the complete author proof (30 body lines).\par\medskip
\noindent\textit{Editorial note.} Navigation stubs: 7 cited-but-not-included labels are supplied as non-load-bearing pointers with their pinned tags where the snapshot tag index has them.\par\medskip
\noindent\textit{Editorial note.} extra\_packages is empty on purpose: the retained preamble already loads xy with options and a second declaration clashes in the assembler wrapper.\par\medskip
\noindent\textit{Editorial note.} Text is retained verbatim, including the author's own typos and placeholders.\par\medskip
\noindent\textit{Editorial note.} Queue index 12 of queue\_pic.json (pb=1539, stmt 624-642, proof 644-675).\par\medskip
\noindent\textit{Editorial note.} stacks-pic wave-1 (pi01): fresh adjudication accept (2026-10-14); pool per standing queue hypercovering -> pic (SMC reassigned to Zhen per Chief 08:5x); pic.tex 27 labels/23 with proof judged 21 accept / 2 reject / 4 pb=0 env-dup rows excluded; spans cite-free and reference/slogan/history-free so no patchenv needed.\par\medskip
\noindent\textit{Allowed standard prerequisites (editorial).}\begin{itemize}
\item Explicit prerequisite boundary. The cited inputs are prerequisites; the argument in the retained span is not replaced by any cited result.
\end{itemize}\medskip
\ifdefined\hypertarget\hypertarget{sourcepointer-1}{}\fi
\noindent\textit{Source reference.} Original lemma 00PD: sources/stacks-pic/pinned-original-source.tex, source lines 675--675. Cited by the retained statement or proof but not part of the collected spans; supplied only as a navigation pointer. The argument does not depend on its content.\par\medskip
\ifdefined\hypertarget\hypertarget{sourcepointer-2}{}\fi
\noindent\textit{Source reference.} Original lemma 0AG8: sources/stacks-pic/pinned-original-source.tex, source lines 675--675. Cited by the retained statement or proof but not part of the collected spans; supplied only as a navigation pointer. The argument does not depend on its content.\par\medskip
\ifdefined\hypertarget\hypertarget{sourcepointer-3}{}\fi
\noindent\textit{Source reference.} Original lemma 062Y: sources/stacks-pic/pinned-original-source.tex, source lines 675--675. Cited by the retained statement or proof but not part of the collected spans; supplied only as a navigation pointer. The argument does not depend on its content.\par\medskip
\ifdefined\hypertarget\hypertarget{sourcepointer-4}{}\fi
\noindent\textit{Source reference.} Original lemma 02LS: sources/stacks-pic/pinned-original-source.tex, source lines 675--675. Cited by the retained statement or proof but not part of the collected spans; supplied only as a navigation pointer. The argument does not depend on its content.\par\medskip
\ifdefined\hypertarget\hypertarget{sourcepointer-5}{}\fi
\noindent\textit{Source reference.} Original lemma 02KB: sources/stacks-pic/pinned-original-source.tex, source lines 675--675. Cited by the retained statement or proof but not part of the collected spans; supplied only as a navigation pointer. The argument does not depend on its content.\par\medskip
\ifdefined\hypertarget\hypertarget{sourcepointer-6}{}\fi
\noindent\textit{Source reference.} Original lemma 01TX: sources/stacks-pic/pinned-original-source.tex, source lines 675--675. Cited by the retained statement or proof but not part of the collected spans; supplied only as a navigation pointer. The argument does not depend on its content.\par\medskip
\ifdefined\hypertarget\hypertarget{sourcepointer-7}{}\fi
\noindent\textit{Source reference.} Original lemma 056S: sources/stacks-pic/pinned-original-source.tex, source lines 675--675. Cited by the retained statement or proof but not part of the collected spans; supplied only as a navigation pointer. The argument does not depend on its content.\par\medskip
\begin{lemma}
\label{lemma-divisors-on-curves}
Let $X \to S$ be a smooth morphism of schemes of relative dimension $1$.
Let $D \subset X$ be a closed subscheme. Consider the following conditions
\begin{enumerate}
\item $D \to S$ is finite locally free,
\item $D$ is a relative effective Cartier divisor on $X/S$,
\item $D \to S$ is locally quasi-finite, flat, and
locally of finite presentation, and
\item $D \to S$ is locally quasi-finite and flat.
\end{enumerate}
We always have the implications
$$
(1) \Rightarrow (2) \Leftrightarrow (3) \Rightarrow (4)
$$
If $S$ is locally Noetherian, then the last arrow is an if and only if.
If $X \to S$ is proper (and $S$ arbitrary), then the first arrow is
an if and only if.
\end{lemma}

\begin{proof}
Equivalence of (2) and (3). This follows from
Divisors, Lemma \ref{divisors-lemma-fibre-Cartier}
if we can show the equivalence of (2) and (3) when
$S$ is the spectrum of a field $k$. Let $x \in X$ be a closed point.
As $X$ is smooth of relative dimension $1$ over $k$ and we see that
$\mathcal{O}_{X, x}$ is a regular local ring of dimension $1$
(see Varieties, Lemma \ref{varieties-lemma-smooth-regular}).
Thus $\mathcal{O}_{X, x}$ is a discrete valuation ring
(Algebra, Lemma \ref{algebra-lemma-characterize-dvr})
and hence a PID. It follows that every sheaf of ideals
$\mathcal{I} \subset \mathcal{O}_X$ which is nonvanishing at all
the generic points of $X$ is invertible
(Divisors, Lemma \ref{divisors-lemma-effective-Cartier-in-points}).
In other words, every closed subscheme of $X$ which does not contain
a generic point is an effective Cartier divisor.
It follows that (2) and (3) are equivalent.

\medskip\noindent
If $S$ is Noetherian, then any locally quasi-finite morphism
$D \to S$ is locally of finite presentation (Morphisms, Lemma
\ref{morphisms-lemma-noetherian-finite-type-finite-presentation}),
whence (3) is equivalent to (4).

\medskip\noindent
If $X \to S$ is proper (and $S$ is arbitrary), then $D \to S$ is
proper as well. Since a proper locally quasi-finite morphism is finite
(More on Morphisms, Lemma \ref{more-morphisms-lemma-characterize-finite})
and a finite, flat, and finitely presented morphism is finite locally free
(Morphisms, Lemma \ref{morphisms-lemma-finite-flat}), we see that
(1) is equivalent to (2).
\end{proof}

\end{document}
